% ============================================
% Johns Hopkins Letter Template
% Assistant Professor Cover Letter – Palm Beach Atlantic University (Rinker College of Business)
% ============================================

\documentclass[11pt,letterpaper]{letter}

\usepackage{graphicx}
\usepackage{eso-pic}
\usepackage{geometry}
\usepackage{microtype}
\usepackage[hidelinks]{hyperref}

% ----- Page Setup -----
\geometry{left=25mm,right=25mm,top=25mm,bottom=25mm}

% ----- Logo Placement (foreground, first page only) -----
\newcommand{\PutJHULogoFG}[1]{%
  \AddToShipoutPictureFG*{%
    \AtPageUpperLeft{%
      \hspace{10mm}%
      \raisebox{-38mm}{%
        \includegraphics[width=6cm]{#1}%
      }%
    }%
  }%
}

% ----- Sender Information -----
\signature{Decory Edwards}
\address{Department of Economics \\ Johns Hopkins University \\ Baltimore, MD 21218 \\ \texttt{dedwar65@jh.edu}}

\begin{document}
\PutJHULogoFG{Images/jhulogo.png}

\begin{letter}{Search Committee \\
Marshall E.~Rinker College of Business \\
Palm Beach Atlantic University \\
901 South Flagler Drive, PO Box 24708 \\
West Palm Beach, FL 33416-4708, USA}

\opening{Dear Members of the Search Committee,}

I am writing to express my interest in the Assistant Professor of Economics position in the Marshall E.~Rinker College of Business at Palm Beach Atlantic University (start date Fall 2026). I am a Ph.D.\ candidate in Economics at Johns Hopkins University, advised by Professors Christopher D.~Carroll, Nicholas W.~Papageorge, and Stelios Fourakis, and I expect to receive my degree in May 2026. My research fields are macroeconomics, wealth and income dynamics, and computational economics.

My research agenda is unified by a focus on understanding the determinants of wealth inequality and the mechanisms through which heterogeneity in returns to wealth shapes aggregate outcomes. In my job market paper, I estimate the degree of heterogeneity in individual portfolio returns using micro data from the Survey of Consumer Finances and simulate a structural model in which households face idiosyncratic return risk. I show that allowing for persistent heterogeneity in returns substantially improves the model’s ability to match the empirical distribution of wealth, offering a tractable way to connect micro-level return dispersion to macroeconomic inequality.

In a second project, I use the Health and Retirement Study to study the relationship between interpersonal trust and portfolio performance. Building on the literature that finds a hump-shaped relationship between trust and income, I show that a similar relationship holds for individual returns to wealth measured in the HRS data, highlighting a behavioral channel through which social attitudes shape saving and investment outcomes. This work also provides new estimates of the persistent component of returns to net worth, which motivate the calibration of heterogeneity in my structural model.

The Rinker College of Business’s AACSB-accredited, mission-driven commitment to equipping students to grow in wisdom and lead with conviction is an excellent fit for my research and methods. My work links household portfolio behavior and return heterogeneity to questions that are central to business decision-making and policy analysis: capital allocation under risk, distributional impacts of taxation, and the interpretation of micro data for forecasting and decision support. I look forward to contributing to undergraduate, MBA, and other graduate programs in a community that integrates faith and learning.

\textbf{Teaching.} My teaching experience centers on undergraduate macroeconomics and an upper-level macro policy seminar, where I have led weekly discussion sections and held regular office hours for classes ranging from 16 to 40 students. Consistent with my teaching statement, my approach is student-centered: I aim to help students ``convince themselves'' that they have moved from confusion to understanding by holding structured office-hour coaching that builds trust and confidence. At PBA, I am prepared to teach economics at the principles, upper-division, and graduate levels, with the flexibility to teach day, evening, on-campus, and online modalities. In addition to principles and intermediate economics, I am able to teach \emph{Principles of Finance} and \emph{Business Statistics}. I would also welcome developing electives such as \emph{Economics of Inequality and Tax Policy} or \emph{Computational Economics for Business Decisions}, emphasizing reproducible workflows and evidence-based decision-making. I embrace PBA’s Christian mission and will thoughtfully integrate faith and ethics into coursework and classroom discussions.

I believe I would be a strong fit because my computational and empirical toolkit allows me to (i) produce robust, data-backed estimates of returns and distributional outcomes, (ii) build and validate models that speak to business, regulatory, and policy-relevant questions aligned with the College’s programs, and (iii) collaborate across teams to present insights in concise, actionable formats for diverse audiences. I am enthusiastic about serving in an AACSB-accredited, Christ-centered business school with a strong emphasis on teaching, scholarship, and service.

I would be honored to contribute to PBA’s research and teaching mission, including a standard three-course load each semester, mentoring students pursuing quantitative careers in economics and business, and engaging in service to the University and wider community. Thank you for your consideration; I welcome the opportunity to discuss my fit with the Rinker College of Business at your convenience.

\closing{Sincerely,}

\end{letter}
\end{document}
